%%%%%%%%%%%%%%%%%%%%%%%%%%%%%%%%%%%%%%%%%%%%%%%%%%%%%%%%%%%%%%%%%%%%%%%%%%%%%%%
\section{Introduction}
\label{sec:intro}

Public radial-velocity (RV) archives now cover two decades. The HARPS-RVBank
\citep{Trifonov2020,Perdelwitz2024} gives uniformly re-derived velocities and
activity indicators for more than 5000 stars observed with HARPS
\citep{Mayor2003}, and it has already been searched for planets
\citep{Sreenivas2022}. Low-amplitude signals in such data have a mixed record.
Signals later withdrawn turned out to be produced by the sampling and the noise
\citep{Rajpaul2016}, by stellar rotation or its harmonics
\citep{Robertson2014,Santos2014,Robertson2015,Carleo2020,Burrows2024}, or by an alias
of rotation that did not hold phase when new data arrived \citep{Lubin2021}.
The tests used to screen candidates are well known individually: false-alarm
probabilities \citep{Baluev2008}, including against red noise \citep{Ejaz2026}, alias comparisons, activity-indicator
periodograms, signal growth and stacked periodograms \citep{MortierCC2017},
apodised sinusoids \citep{Hara2022b} and false-inclusion probabilities
\citep{Hara2022}. What is rarely measured is how well each of them separates
planets from impostors on the sampling and noise of real archival series, and
how they perform on signals whose fate is known.

Blind challenges have measured how well activity models recover planets in
simulated data sets \citep{Dumusque2017,Zhao2022}, and a neural network trained on
simulations has been proposed to replace the significance test and classify signals as
planetary or not \citep{Nieto2023}. Here we ask a complementary
question for archival searches: given a significant periodic signal in a long,
irregular, single-instrument series, which tests tell a planet from an
impostor, and how often do they fail? We answer it in four steps.
\begin{enumerate}
\item We implement a ladder of interpretable tests as tested, open software
      (Sect.~\ref{sec:ladder}).
\item We calibrate every test on simulations built from the real sampling and
      noise of \nmSimStars\ HARPS-RVBank stars, with planets and evolving stellar
      activity injected and every significant peak labelled
      (Sects.~\ref{sec:sims} and~\ref{sec:simres}).
\item We apply the ladder, a rule whose thresholds were fixed before the benchmark was run, and
      classifiers trained only on the simulations to published signals around
      HARPS-RVBank stars whose later verdict is known (Sect.~\ref{sec:bench}), and to the
      archive as it stood in January 2022, against the planets published since
      (Sect.~\ref{sec:timesplit}).
\item Because the analysis that motivated this work was written with an AI
      assistant, we report the defects that checks of it found, coded by two
      independent instances of the model, as a case study in verifying AI-assisted analysis
      (Sect.~\ref{sec:verify}).
\end{enumerate}
The candidate that motivated the work, HD\,297396\,b (Fraser, in preparation), is
the worked example of Sect.~\ref{sec:example}.

%%%%%%%%%%%%%%%%%%%%%%%%%%%%%%%%%%%%%%%%%%%%%%%%%%%%%%%%%%%%%%%%%%%%%%%%%%%%%%%
\section{Data}
\label{sec:data}

We use the HARPS-RVBank table of \citet{Perdelwitz2024} (CDS J/A+A/683/A125, the
version of 2026 April 20), with \nmRVBankSpectra\ spectra of \nmRVBankStars\ stars
taken before 2022. For each spectrum it gives the \texttt{SERVAL} velocity
\citep{Zechmeister2018} corrected for the nightly zero point, its uncertainty, and
eight activity indicators: the differential line width (dLW), H$\alpha$, the
Na\,\textsc{i} D1 index, the chromatic index (CRX), the full width at half maximum
(FWHM), contrast and bisector span of the cross-correlation function, and
$R'_{\rm HK}$.

The same instrumental cuts, none of which uses model residuals, apply to every
star: RVBank flag zero, signal-to-noise ratio at least 10 in order 55, absolute
drift at most $3\ms$, and a finite uncertainty below $20\ms$ (\nmRVBankKeptPct\% of
all spectra pass). Spectra are averaged by night with inverse-variance weights,
separately before and after the 2015 fibre upgrade \citep{LoCurto2015}; a night runs
from local noon to local noon at La Silla (16:43\,UT). The two eras receive
independent offsets and jitters. For every star, simulated or real, nightly epochs
more than 15 robust standard deviations from a per-era quadratic are dropped, which
removes misidentified or grossly deviant spectra (\nmBenchClipped\ epochs in
\nmBenchClippedStars\ benchmark hosts) but keeps merely discrepant nights, and a
quadratic trend is fitted jointly with everything else. We write $N$ for the number of
nightly epochs, $T$ for the baseline and $n_{\rm era} \le 2$ for the number of eras.

%%%%%%%%%%%%%%%%%%%%%%%%%%%%%%%%%%%%%%%%%%%%%%%%%%%%%%%%%%%%%%%%%%%%%%%%%%%%%%%
\section{The vetting ladder}
\label{sec:ladder}

\subsection{Periodogram and detection}
\label{sec:detect}

For velocities $y_i$ with uncertainties $\epsilon_i$ and era labels $g(i)$, the model
has per-era offsets and per-era jitters $s_g$, so that the variance of epoch $i$ is
$\sigma_i^2 = \epsilon_i^2 + s_{g(i)}^2$. With the jitters fixed, the improvement in
$\chi^2$ from adding a sinusoid at frequency $f$ is
\begin{equation}
\dchi(f) = \chi^2_{\rm offsets} - \chi^2_{\rm offsets + a\cos 2\pi ft + b \sin 2\pi ft},
\qquad \chi^2 = \sum_i r_i^2/\sigma_i^2 .
\end{equation}
Offsets are removed by projecting the whitened sinusoids and data (each divided by
$\sigma_i$) onto the complement of the whitened era indicators, so each frequency
costs $O(N n_{\rm era})$ operations and many simulated series are processed in one
matrix product. The search band is 1.2--500\,d on a grid of three frequencies per
$1/T$.

The jitters matter. We first fit them by maximum likelihood with offsets only, and
take the candidate frequency $f_0$ as the highest peak of $\dchi(f)$ with those
jitters (or, for a claimed period, the highest point within $\pm0.5/T$ of it, half a
resolution element on either side; Sect.~\ref{sec:changes} gives the effect of this
choice). A real
signal inflates the planet-free jitters and so deflates its own $\dchi$. We therefore
rescale all jitters by one common factor fitted with the sinusoid at $f_0$ in the
model, and use as detection statistic the likelihood ratio
$z_0 = 2(\ln L_1 - \ln L_0)$, with the jitters at their planet-free values under $L_0$
and rescaled under $L_1$. (Refitting each era's jitter separately under $L_1$ can
collapse the jitter of an era with few epochs onto a sinusoid that happens to fit
them; one common factor cannot.) All later rungs use the rescaled jitters. The
false-alarm probability (FAP) of $z_0$ over the band is the approximation of
\citet{Baluev2008} for a sinusoid in noise of known variance,
\begin{equation}
{\rm FAP} \simeq 1 - (1 - e^{-z_0/2})\,e^{-\tau}, \quad
\tau = W e^{-z_0/2}\sqrt{z_0/2},
\end{equation}
with $W = (f_{\max} - f_{\min})\sqrt{4\pi\,{\rm Var}_w(t)}$ and ${\rm Var}_w(t)$ the
variance of the epochs weighted by $1/\sigma_i^2$. It bounds the FAP from above as the
FAP tends to zero; how it performs on real noise with a fitted jitter is measured in
Sect.~\ref{sec:fap}. For a period known in advance, the single-frequency probability
is $e^{-z_0/2}$.

\subsection{The rungs}

Given $f_0$, the ladder computes the features of Table~\ref{tab:features}.

\emph{Amplitude.} The semi-amplitude $K$ and its uncertainty from the sinusoid fit,
$K/\sigma_K$, and $K$ over the rms of the velocities about the offsets.

\emph{Period.} The candidate aliases of $f_0$ are its first-order aliases through
the sidereal day, the synodic month and the year and through the three strongest
peaks of the star's own spectral window. Those inside the search band are refined
within $\pm 0.5/T$ (inside the band), and any that refines to within $1.5/T$ of $f_0$ is
dropped, since it has found the candidate's own peak. The feature is the margin in
$\dchi$ between $f_0$ and the strongest remaining alias.

\emph{Stationarity.} A Keplerian keeps amplitude and phase; rotational modulation
does not. We fit the sinusoid at $f_0$, with a common phase origin, separately to the
two halves of the series and to four blocks with equal numbers of epochs. Each fit
gives $(a, b)$ with a covariance matrix. About their generalised-least-squares mean
$\mu$, the deviations are split into the component along $\mu$ (amplitude) and across
it (phase), and each sum of squared standardised deviations is compared with a
$\chi^2$ distribution with $n_{\rm b} - 1$ degrees of freedom, $n_{\rm b}$ being the
number of sub-samples. For four blocks this gives the probabilities
\code{block\_phase\_p} and \code{block\_amp\_p}, and the rms phase deviation; for
the halves, the phase and amplitude differences in units of their uncertainty and the
imbalance between the halves' share of $\dchi$ and their share of epochs. Because $f_0$
is fitted to the whole series, any constant phase drift between the halves is absorbed
into it, so the half-phase test is weak by construction; the block test, with four
sub-samples, retains power, and its calibration on simulated planets is given in
Sect.~\ref{sec:calib}. We also follow $\dchi$ as epochs accumulate in time order and
report its Spearman correlation with the number of epochs
\citep[cf.][]{MortierCC2017,Sreenivas2022}, and we replace the constant amplitude by a
Gaussian envelope over a grid of centres and widths (0.1--0.5 of the baseline) and
report the gain in $\dchi$ \citep[cf. the apodised sinusoids of][]{Hara2022b}.

\emph{Activity.} Each indicator is fitted with per-era offsets and the same quadratic trend
as the velocities, which are removed, and is divided by its robust scatter over the whole
series. Values more than 15 robust standard deviations from the fit, the level used for the
velocities, are dropped and the fit repeated, since a least-squares trend tilted by one
gross outlier leaves a spurious long-period signal. Indicators with values on fewer than
60\% of the epochs are dropped. Indicators carry no usable errors, so each era receives the
closed-form maximum-likelihood jitter (the era's standard deviation). The feature
\code{ind\_max\_dchi2} is the largest $\dchi$ of any indicator within $\pm 0.5/T$ of
$f_0$ (eleven frequencies); with no usable indicator it is zero. We also report the
largest absolute correlation of an indicator with the velocities. As a rotation proxy
we take, among dLW, H$\alpha$, FWHM and $R'_{\rm HK}$, the most significant local maximum
over 2.5--200\,d, excluding frequencies within $1/T$ of one and two cycles per year, where
seasonal sampling puts spurious power; its Baluev FAP is multiplied by the number of
indicators searched, and we write $f_{\rm proxy}$ for its frequency. When that FAP is below 1\%, \code{harm\_dist} is
the distance of $f_0$ from the nearest of $k f_{\rm proxy}$, $k = 1$--4, in units of $1/T$, and \code{harm\_alias\_dist} the
distance from the nearest first-order alias of one of them; otherwise both are
undefined and the classifiers receive a flag, \code{harm\_missing}.

\emph{Known signals.} When a system has accepted planets, they are removed first by
Keplerian fits with fixed periods. For fixed eccentricity and time of periastron the
orbit is linear in $K\cos\omega$ and $K\sin\omega$, so the orbital elements are found by
a grid in $(e, t_p)$ and simplex refinement, by backfitting over the planets, jointly
with the offsets and the quadratic trend.

\begin{table}
\caption{Ladder features and the question each answers.}
\label{tab:features}
\centering
\footnotesize
\setlength{\tabcolsep}{3pt}
\begin{tabular}{L{2.6cm}L{5.4cm}}
\hline\hline
Feature & Definition \\
\hline
\code{dchi2}, \code{log10\_fap} & likelihood ratio $z_0$; its Baluev FAP over 1.2--500\,d \\
\code{snr\_K}, \code{K\_over\_rms} & $K/\sigma_K$; $K$ over the rms about the offsets \\
\code{alias\_margin} & $\dchi(f_0)$ minus the largest alias $\dchi$ inside the band \\
\code{half\_phase\_z}, \code{half\_amp\_z} & phase, amplitude difference between halves (in $\sigma$) \\
\code{half\_imbalance} & $|$share of $\dchi$ $-$ share of epochs$|$ in the first half \\
\code{block\_phase\_p}, \code{block\_amp\_p}, \code{block\_phase\_rms} & constancy of phase and amplitude in four blocks; rms phase deviation \\
\code{growth\_rho} & Spearman $\rho$ of $\dchi(f_0)$ with epochs in time order \\
\code{apod\_gain} & $\dchi$ gain of a Gaussian envelope \\
\code{ind\_max\_dchi2}, \code{ind\_max\_absr} & largest indicator $\dchi$ near $f_0$; largest $|r|$ with the RVs \\
\code{act\_log10\_fap} & FAP of the rotation proxy (2.5--200\,d, seasonal aliases excluded) \\
\code{harm\_dist}, \code{harm\_alias\_dist}, \code{harm\_missing} & distance to $k f_{\rm proxy}$ and to its aliases ($1/T$ units); no proxy \\
\code{log\_n}, \code{log\_ncyc} & $\log N$; $\log (T f_0)$. Reported, not used by the classifiers \\
\hline
\end{tabular}
\end{table}

\subsection{An a-priori rule}
\label{sec:rule}

Before any benchmark signal was run through the pipeline we wrote down a rule-based
verdict (released with the code): a candidate passes if its FAP is below 1\%, it
beats its aliases, its phase is constant across blocks ($p > 0.01$) and halves
($<3\sigma$), no indicator has $\dchi \ge 13.8$ near $f_0$, and it is not within one
resolution element of a rotation harmonic. A rung that cannot be computed does not
reject: with fewer than three usable blocks or no halves the phase rung passes, and
without a significant rotation proxy the rotation rung passes. The thresholds came from nominal
significance levels and from the analysis of HD\,297396\,b, so that candidate's
passing the rule (Sect.~\ref{sec:example}) is partly circular. They are not nominal
error rates: 13.8 is the single-frequency $p = 10^{-3}$ point of a $\chi^2_2$
distribution, but it is applied to the maximum over eight indicators and eleven
frequencies. We therefore report what each rung does to simulated planets
(Sect.~\ref{sec:calib}). The rule does not use \code{harm\_alias\_dist}, which we
added after the first benchmark pass; results with it are labelled as such. The
thresholds have not changed since they were frozen, but the features they apply to
have (Sect.~\ref{sec:changes}).

%%%%%%%%%%%%%%%%%%%%%%%%%%%%%%%%%%%%%%%%%%%%%%%%%%%%%%%%%%%%%%%%%%%%%%%%%%%%%%%
\section{Simulations on real sampling and noise}
\label{sec:sims}

We selected \nmSimStars\ stars at random from those with 40--250 nightly epochs, a
baseline of at least 1500\,d, a median uncertainty of at most $3\ms$, dLW and H$\alpha$
on at least 80\% of the spectra, and an RV standard deviation of at most $20\ms$
(which excludes binaries and giant planets), and excluded every benchmark host and
HD\,297396.

\emph{Noise.} Each simulated series starts from the star's real nightly velocities
and indicators with one random permutation within each era, shared by the
velocities, the indicators and the uncertainties. The sampling, the errors and the
joint distribution of velocity and indicator noise, including its tails, are real;
any coherent signal the star had, and any noise correlated between nights, is
destroyed. The star's series is clipped for outliers before the permutation, and each
simulated series is fitted with the same quadratic trend as a real star.

\emph{Planets.} Keplerians with $P$ log-uniform in 1.2--300\,d, eccentricity from
the distribution of \citet{Kipping2013} redrawn above 0.8, random orientation and
phase, and $K$ log-uniform in 0.15--3 times the scatter $\sigma$ of the noise series.
We describe each case by its expected $K/\sigma_K$, with
$\sigma_K = \sigma\sqrt{2/N}$.

\emph{Activity.} We use the FF$'$ picture \citep{Aigrain2012,Rajpaul2015}: a latent
process $G(t)$ with a quasi-periodic kernel
$k(\tau) = \exp[-\tau^2/2\lambda^2 - \sin^2(\pi\tau/P_{\rm rot})/2w^2]$ and its
derivative $G'(t)$, drawn jointly at the epochs from the analytic covariance of
$(G, G')$. The velocities receive $A[\sqrt{1-\varphi^2}\,G + \varphi G']$, $\varphi$ being
the share of the rotational term $G'$, with
$P_{\rm rot}$ log-uniform in 3--150\,d, $\lambda$ log-uniform in 1--30 rotations, $w$
uniform in 0.25--1.2, $A$ log-uniform in 0.2--3 times $\sigma$ and $\varphi$ uniform in 0--0.9. Both $K$ and $A$ scale with $\sigma$ and not with $N$, so
the number of epochs carries no information about the class. How strongly activity
shows in the indicators is the most consequential assumption of the simulations, so
we vary it: in four cases out of five each indicator $j$ receives
$A_j[\cos\theta_j\,G + \sin\theta_j\,G']$ times its own scatter, with a common amplitude
log-uniform in 0.1--3, a factor $10^{\mathcal{N}(0,\,0.3)}$ for each indicator and a
phase lag $\theta_j$ uniform within $\pm45\degr$ of $0$ (all indicators except the
bisector) or of $90\degr$ (bisector); in the fifth case the indicators receive nothing.
The indicator amplitudes are drawn independently of the RV amplitude $A$.

\emph{Cases and labels.} Each case is a planet alone (35\%), a planet with activity
(25\%), activity alone (30\%) or noise alone (10\%). The ladder is run blind on the
highest peak. That peak is labelled, in this order of precedence:
a \emph{planet} if it lies within $1/T$ of the injected orbital frequency $f_p$;
a \emph{planet alias} if within $1/T$ of $2f_p$, $3f_p$, $f_p/2$ or a first-order alias
of $f_p$; a \emph{rotation harmonic} if within $1/T$ of $k f_{\rm rot}$ ($k = 1$--4, with
$f_{\rm rot} = 1/P_{\rm rot}$ the injected rotation frequency) or a first-order alias of
one; a \emph{rotation line} if within
$\max[1/T, 1/(\pi\lambda)]$ of $k f_{\rm rot}$, $k = 0$--4, that is inside the
$\pm2\sigma$ width of a line of the quasi-periodic power spectrum, whose lines are
Gaussians of standard deviation $1/(2\pi\lambda)$ (for $k = 0$, inside the
low-frequency envelope of the evolving amplitude); otherwise \emph{other activity} if
the series has activity, and \emph{noise} if not. Every case is reproducible from its
seed.
